\section{Quantum (Bayesian) Proof-Nets}
\label{sec:qpns}

In this section, we introduce \qpns.
They can be thought of as a typed version of \QBN s, built on the graph formalism of Multiplicative Linear Logic.
Any \QBN\ can be encoded as a \qpn\ and, conversely, to any \qpn of appropriate type is associated a \QBN.
What we gain by moving to the typed setting of proof-nets is:
\begin{itemize}
\item 
This formalism encodes \qqbn s and enables formalizing (and working with) sub-networks.
\item
Typing allows us to \emph{modularly compose} \qpns of compatible interfaces.
\item
Finally, we have \emph{a linear form of high-order}, as already sketched in \cref{fig:modular} where the network $\N_1$ of type $(A\lollipop C) \lollipop D$ is expecting as input a network of type $(A\lollipop C)$.
\end{itemize}

\subparagraph{Formulas.}

We assume given two countable sets of \textdef{names} denoted by metavariables $X, Y, \dots$ and $q_1, q_2, \dots$ which we call respectively \textdef{classical names} and \textdef{q-names}.
Formulas are those of the Multiplicative fragment of Linear Logic (\MLL):
\begin{equation*}
	\FormA,\FormB ::=
	X^+	\mid X^- \mid q^+ \mid q^-
	\mid\FormA\otimes\FormB
	\mid\FormA\parr\FormB
\end{equation*}

\textdef{Negation} $(\cdot)\orth$ is defined inductively by
$(X^+)\orth\defeq X^-$,
$(X^-)\orth\defeq X^+$, 
$(q^+)\orth \defeq q^-$,
$(q^-)\orth \defeq {q^+}$,
$(\FormA\otimes\FormB)\orth \defeq\FormA\orth\parr\FormB\orth$ and
$(\FormA\parr\FormB)\orth \defeq\FormA\orth\otimes\FormB\orth$.
\textdef{Linear arrow} is defined as usual by $F\lollipop G \defeq F^\bot \parr G$. 
We reserve the notation $\qreg^+$ for formulas of the shape $\otimes_{i\in I} q_i^+$.
Capital Greek letters $\Gamma,\Delta,\dots$ vary over finite sequences of formulas, \ie\ $\Delta = \FormA_1,\dots, \FormA_n$.
By $\Nm F$ we denote the set of all names which appear in the formula $F$, and similarly for $\Nm \Delta$.
For example, $\Nm{X^-\parr (Y^+\otimes X^+), q^+\parr q^-} =\{X,Y,q\}$.

A key property of formulas in Linear Logic is (positive/negative) \textdef{polarity}, which we indicate for atoms and extend to formulas as follows, defining (strictly) polarized formulas:
\begin{align*}
\textbf{Positive formulas: }& P, P'::= X^+ \mid q^+ \mid P \otimes P'
\\
\textbf{Negative formulas: }& N, N'::= X^- \mid q^- \mid N \parr N'
\end{align*}

\begin{remark}\label{rem:polarity}
Dual polarities carry a connotation of output/input, answer/question, active/ passive, player/opponent, etc.
Negative atoms are seen as \emph{inputs}, and positive ones as \emph{outputs}.
An intuition---at the base of the Linear Logic interpretation of the information flow---is that information travels upwards (\resp\ downwards) on negative (\resp\ positive) atoms.
\end{remark}

\subparagraph{MLL proof-nets (with boxes).}

In Linear Logic, proofs admit both a \emph{sequent calculus} syntax in the form of \textdef{\pts} and a \emph{graph syntax}, in the form of \textdef{\pn}.
Here, we are interested in the latter.
We omit the sequent calculus presentation---which is straightforward to deduce for the reader familiar with Linear Logic, but not relevant for our purpose.

\begin{definition}[Proof-Nets]
\label{def:pn}
We call \textbf{typed} a labelled (partial\footnote{A partial graph may have \emph{pending edges}, which may be either out-edges (as in \cref{ex:bell_pn}) or in-edges.}) graph $\somepn$ built from the alphabet of nodes on \cref{fig:grammar_nodes_pn}, with edges labelled by \MLL formulas.
Edges incident to a node $\somenode$ need to respect the typing conditions of the grammar, and are classified either as \definitive{premises} of $\somenode$ (depicted above $\somenode$) or as its \definitive{conclusions} (depicted below $\somenode$).
We further require that an edge is the conclusion of at most one node and the premise of at most one node.
\begin{itemize}
\item
\textbf{Proof-nets.}
A typed graph $\somepn$ is a \definitive{proof-net} if it has no pending premises (\ie\ every edge in $\somepn$ is the conclusion of some node) and if it is \definitive{correct}:
every cycle in $\somepn$ uses the two premises of a same $\parr$-node or the two premises of a same $\cc$-node.
\item
\textbf{Conclusions.}
The labels of the edges which are premise of no node (\aka\ pending conclusions) are called the \definitive{conclusions} of $\somepn$.
The \definitive{type} of a \pn $\somepn$ of conclusions $\Delta=C_1, \dots, C_k$ is the sequent $\vdash \Delta$.
We also write $\somepn \der \Delta$.
\end{itemize}
\end{definition}

\begin{example}
The typed graph in \cref{ex:bell_pn}---where $\qreg^+=q_1^+\tens q_2^+$---is a \pn.
\end{example}

The grammar in \cref{fig:pn} extends the standard nodes of \MLL\ with a new node, called a \textbf{box}, which has \emph{exactly one positive conclusion}---its \definitive{\main}---of type either $\X$ or $Q^+=\otimes_{i\in I} q_i^+$.
The negative conclusions---called \definitive{inputs}---are atomic.
All atoms appearing in the conclusions of a box are pairwise distinct.
The grammar is similar to that introduced in~\cite{EhrhardFP23}.
However, while in~\cite{EhrhardFP23} to each box is associated a CPT, here we associate a \qCPT.

\begin{figure}
\begin{adjustbox}{}
% Ax
\begin{tikzpicture}
\begin{scope}[every node/.style={draw=none}, every path/.style={draw=black}]
	\node (ax) at (0,0) {$\ax$};
	\coordinate (axc1) at (-.75,-.5);
	\coordinate (axc2) at (.75,-.5);
	\path (axc1) |- node[named edge,pos=.25]{$X^-$} (ax) -| node[named edge,pos=.75]{$X^+$} (axc2);
\end{scope}
\end{tikzpicture}
\enskip
\begin{tikzpicture}
\begin{scope}[every node/.style={draw=none}, every path/.style={draw=black}]
	\node (ax) at (0,0) {$\ax$};
	\coordinate (axc1) at (-.75,-.5);
	\coordinate (axc2) at (.75,-.5);
	\path (axc1) |- node[named edge,pos=.25]{$q^-$} (ax) -| node[named edge,pos=.75]{$q^+$} (axc2);
\end{scope}
\end{tikzpicture}
\enskip
% Cut
\begin{tikzpicture}
\begin{scope}[every node/.style={draw=none}, every path/.style={draw=black}]
	\node (ax) at (0,0) {$\cut$};
	\coordinate (axc1) at (-.75,.5);
	\coordinate (axc2) at (.75,.5);
	\path (axc1) |- node[named edge,pos=.25,left]{$A\orth$} (ax) -| node[named edge,pos=.75,right]{$A$} (axc2);
\end{scope}
\end{tikzpicture}
\enskip
% ⊗
\begin{tikzpicture}
\begin{scope}[every node/.style={draw=none}, every path/.style={draw=black}]
	\node (c) at (0,0) {$\tens$};
	\coordinate (cp1) at (-.5,.5);
	\coordinate (cp2) at (.5,.5);
	\coordinate (cc1) at (0,-.75);
	\path[out=-90,in=180] (cp1) edge node[named edge,left]{$A$} (c);
	\path[out=-90,in=0] (cp2) edge node[named edge,right]{$B$} (c);
	\path (c) -- node[named edge]{$A\tens B$} (cc1);
\end{scope}
\end{tikzpicture}
\enskip
% `
\begin{tikzpicture}
\begin{scope}[every node/.style={draw=none}, every path/.style={draw=black}]
	\node (c) at (0,0) {$\parr$};
	\coordinate (cp1) at (-.5,.5);
	\coordinate (cp2) at (.5,.5);
	\coordinate (cc1) at (0,-.75);
	\path[out=-90,in=180] (cp1) edge node[named edge,left]{$A$} (c);
	\path[out=-90,in=0] (cp2) edge node[named edge,right]{$B$} (c);
	\path (c) -- node[named edge]{$A\parr B$} (cc1);
\end{scope}
\end{tikzpicture}
\enskip
% Contraction
\begin{tikzpicture}
\begin{scope}[every node/.style={draw=none}, every path/.style={draw=black}]
	\node (c) at (0,0) {$\cc$};
	\coordinate (cp1) at (-.5,.5);
	\coordinate (cp2) at (.5,.5);
	\coordinate (cc1) at (0,-.75);
	\path[out=-90,in=180] (cp1) edge node[named edge,left]{$\negatom{X}$} (c);
	\path[out=-90,in=0] (cp2) edge node[named edge,right]{$\negatom{X}$} (c);
	\path (c) -- node[named edge]{$\negatom{X}$} (cc1);
\end{scope}
\end{tikzpicture}
\enskip
% Weakening
\begin{tikzpicture}
\begin{scope}[every node/.style={draw=none}, every path/.style={draw=black}]
	\node (w) at (0,0) {$\cw$};
	\coordinate (cc1) at (0,-.75);
	\path (w) -- node[named edge]{$X^-$} (cc1);
\end{scope}
\end{tikzpicture}
\enskip
\begin{tikzpicture}
\begin{scope}[every node/.style={draw=none}, every path/.style={draw=black}]
	\node (w) at (0,0) {$\cw$};
	\coordinate (cc1) at (0,-.75);
	\path (w) -- node[named edge]{$q^-$} (cc1);
\end{scope}
\end{tikzpicture}
\enskip
% boxnc
\begin{tikzpicture}
\begin{scope}[every node/.style={rectangle,draw=black}, every path/.style={draw=black}]
	\node[minimum width=25mm,minimum height=6mm] (b) at (0,0) {$X/\qreg$};
	\coordinate (bc1) at (-1,-1);
	\coordinate (bc2) at (-.2,-1);
	\coordinate (bcn) at (1,-1);
	\node[draw=none] at (-.6,-.5) {$\dots$};
	\node[draw=none] at (.5,-.5) {$\dots$};
	\path (-1,-.3) -- node[named edge,left,pos=.75]{$\negatom{Y_1}$} (bc1);
	\path (-.2,-.3) -- node[named edge,pos=.75]{$\negatom{q_1}$} (bc2);
	\path (1,-.3) -- node[named edge,right,pos=.75]{$\posatom{X}/\qreg^+$} (bcn);
\end{scope}
\end{tikzpicture}
\end{adjustbox}
\caption{
	Grammar of nodes for proof-nets;
%	For a $\cw$-node, $N$ is a negative formula.
	atoms in the conclusions of a box are pairwise distinct.
}
\label{fig:grammar_nodes_pn}\label{fig:pn}
\end{figure}

\subparagraph{Quantum \pns.}

As before, we identify each name $X$ with the \rv\ $(X, \{x^{\true}, x^{\false}\})$ and each q-name $q_i$ with a distinct qubit.
With this implicit assumption, if $P$ is the \main\ of a box $\bbox$, then $\Nm P$ is either a random variable $\{X\}$ or a quantum register $\{q_1, \dots, q_k\}$.

\begin{definition}[Quantum \pns]\label{def:qbn}
A \definitive{\qpn} (\qbpn) is a pair $(\somepn, \Phi)$ where:
\begin{enumerate}
\item\label{item:qbn:1}
$\somepn$ is a proof-net such that all the atoms appearing in the outputs of distinct boxes are pairwise distinct, and all the negative q-atoms which are inputs of distinct boxes are pairwise distinct.
\item
$\someqfactset$ assigns to each box $\bbox$ of \main $P$ a \qCPT\ $\someqfactset(\bbox)$, for $\Nm{P}$ given the (classical and q-) names of the inputs of $\bbox$.
\end{enumerate}
\end{definition}

Thanks to point \ref{item:qbn:1}, we denote a box in a \qbpn $\somepn$ by its positive conclusion---\eg\ $\boxof{P}$ for the unique box of \main $P$. When $\Phi$ is irrelevant, we often write ``a \qbpn $\somepn$'' for ``a \qbpn $(\somepn,\Phi)$''.

\subparagraph{Reduction rules \& Normal forms.}

\Qpns are equipped with the standard reduction rules for \MLL proof-nets, presented in \Cref{fig:rewriting_pn}.
These rules define a binary relation $\Red$ on \qbpns, written $\somepn \Red \somepn'$ (read $\somepn$ reduces to $\somepn'$).
A \qbpn $\somepn$ is in \definitive{normal} form (or just normal) if there is no $\somepn'$ such that $\somepn \Red \somepn'$. 
Please notice that, because of boxes, a proof-net in normal form can still contain $\cut$-nodes---see \eg\ the \pn in \cref{ex:bell_pn}.

\begin{figure}
\begin{adjustbox}{}
% Ax-cut
\begin{tikzpicture}[baseline=(base)]
\coordinate (base) at (0,-.25);
\begin{scope}[every node/.style={draw=none}, every path/.style={draw=black}]
	\node (ax) at (0,0) {$\ax$};
	\node (cut) at (1.5,-.5) {$\cut$};
	\coordinate (c) at (-.75,-.5);
	\coordinate (m) at (.75,-.25);
	\coordinate (h) at (2.25,0);
	\path (ax) -| node[named edge,pos=.75,left]{$A\orth$} (c);
	\path (ax) -| node[named edge,pos=1]{$A$} (m);
	\path (cut) -| (m);
	\path (cut) -| node[named edge,pos=.75,right]{$A\orth$} (h);
\end{scope}
\end{tikzpicture}
$\quad\Red\quad$
\begin{tikzpicture}[baseline=(base)]
\coordinate (base) at (0,0);
\begin{scope}[every node/.style={draw=none}, every path/.style={draw=black}]
	\coordinate (c) at (0,-.5);
	\coordinate (h) at (0,.5);
	\path (c) -- node[named edge,pos=.5,right]{$A\orth$} (h);
\end{scope}
\end{tikzpicture}
\qquad\vrule\qquad
% ⊗\`
\begin{tikzpicture}[baseline=(base)]
\coordinate (base) at (0,0);
\begin{scope}[every node/.style={draw=none}, every path/.style={draw=black}]
	\node (tens) at (-.75,0) {$\tens$};
	\coordinate (tens1) at (-1.25,.5);
	\coordinate (tens2) at (-.25,.5);
	\node (par) at (1.75,0) {$\parr$};
	\coordinate (par1) at (1.25,.5);
	\coordinate (par2) at (2.25,.5);
	\node (cut) at (.5,-.5) {$\cut$};
	\path[out=-90,in=180] (tens1) edge node[named edge,left]{$A$} (tens);
	\path[out=-90,in=0] (tens2) edge node[named edge,right]{$B$} (tens);
	\path (tens) |- node[named edge,left,pos=.3]{$A \tens B$} (cut);
	\path[out=-90,in=180] (par1) edge node[named edge,left]{$A\orth$} (par);
	\path[out=-90,in=0] (par2) edge node[named edge,right]{$B\orth$} (par);
	\path (par) |- node[named edge,right,pos=.3]{$A\orth \parr B\orth$} (cut);
\end{scope}
\end{tikzpicture}
$\quad\Red\quad$
\begin{tikzpicture}[baseline=(base)]
\coordinate (base) at (0,0);
\begin{scope}[every node/.style={draw=none}, every path/.style={draw=black}]
	\coordinate (tens1) at (-1.25,.5);
	\coordinate (tens2) at (-.25,.5);
	\coordinate (par1) at (1.25,.5);
	\coordinate (par2) at (2.25,.5);
	\node (cut1) at (0,0) {$\cut$};
	\node (cut2) at (1,-.5) {$\cut$};
	\path[out=-90,in=180] (tens1) |- node[named edge,left,pos=.3]{$A$} (cut1);
	\path[out=-90,in=0] (tens2) |- node[named edge,left,pos=.4]{$B$} (cut2);
	\path[out=-90,in=180] (par1) |- node[named edge,right,pos=.3]{$A\orth$} (cut1);
	\path[out=-90,in=0] (par2) |- node[named edge,right,pos=.4]{$B\orth$} (cut2);
\end{scope}
\end{tikzpicture}
\end{adjustbox}
\caption{Reduction rules of proof-nets.}
\label{fig:rewriting_pn}
\end{figure}

\begin{proposition}
The reduction $\Red$ preserves both the correctness and the conclusions of a proof-net, is terminating and is confluent.
\end{proposition}

\subsection{Quantum Bayesian Networks and Quantum Proof-Nets}\label{sec:QBN_qpn}

We now relate the two formalisms, quantum Bayesian networks and quantum proof-nets.

\subparagraph{Polarized proof-nets as directed graphs.}

Polarities embed a notion of orientation (\Cref{rem:polarity}), allowing to read a typed graph as a DAG.
Restricting our attention to proof-nets whose edges are labelled only by polarized formulas, known as \textdef{polarized proof-nets}~\cite{Laurent03}, to each edge is associated a \emph{direction}: upwards for edges labelled by a negative formula, downwards for edges labelled by a positive formula.
The following is well-known for polarized proof-nets.

\begin{lemma}[Polarized correctness~\cite{phdlaurent}]\label{lem:pol_correct}
Let $\somepn$ be a typed graph (\Cref{def:pn}) whose edges are labelled only by polarized formulas. 
The graph $\somepn$ is \emph{correct} if and only if orienting the edges of $\somepn$ according to their polarity yields a DAG.
\end{lemma}

By polarized correctness, every polarized \pn\ $\somepn$ can be seen as a DAG where edges are oriented according to their polarity. 
This induces a DAG on the boxes $\Boxes\somepn=\{\boxof{P_1}, \dots, \boxof{P_n}\}$ of $\somepn$, hence a DAG $\G_\somepn$ on $\Nm{P_1}, \dots, \Nm{P_n}$: we have a directed edge $\Nm{P_i} \rightarrow \Nm{P_j}$ in $\G_\somepn$ whenever there is a directed path from $\bbox^{P_i}$ to $\bbox^{P_j}$ in $\somepn$.
%If we have a \qbpn $(\R,\Phi)$, it is straightforward to equip $\G_\R$ with a function associating a \qCPT to each node, in the same way as $\Phi$ does.

\subparagraph{From \qpns to \qqbn's.}

%The former formalism is a more expressive syntax, so we relate \qqbn's with the class of \qpns with ``no inputs''.
We call \definitive{\good} a \qpn $\somepn \der \Delta$ where all the atoms appearing (as sub-formulas) in $\Delta$ are \emph{positive and classical}.

\begin{lemma}
To each \emph{polarized} \qbpn $(\somepn,\Phi_\somepn)$ of boxes $\Boxes{\somepn}=\{\bbox^{P_1}, \dots, \bbox^{P_n}\}$ is associated a pair $(\G, \Phi_\G)$ with $\G$ the DAG induced by $\somepn$ on $\Nm{P_1}, \dots, \Nm{P_n}$ and $\Phi_\G(\Nm P)\defeq \Phi_\somepn(\bbox^P)$.
If $\somepn$ is \good, the pair $(\G, \Phi_\G)$ is a \QBN.
\end{lemma}
 
\begin{example}
The \pn\ in \cref{ex:bell_pn} is polarized and \good; its DAG is the one in \Cref{fig:Bell}.
\end{example}

% Notice that the \BN associated to a \BPN $\R$ is independent from the conclusions of $\R$.

\subparagraph{From \qqbn s to \qpns.}

A quantum Bayesian network can be encoded as a \qpn rather directly.
%\footnote{Similarly to the encoding of \BNs into \pns~\cite{old_BPN}.}
To a \qqbn\ $(\G, \Phi)$, it is straightforward to associate \emph{a typed graph} $\Net_\G$ as sketched in \cref{fig:Bell,ex:bell_pn}.
It is immediate that $\Net_\G$ is polarized and \good.
The only delicate point is checking $\Net_\G$ is correct, hence a proof-net: it follows from \cref{lem:pol_correct}, as $\G$ is a DAG.
The pair $(\Net_\G, \Phi)$ clearly satisfies all conditions in \cref{def:qbn}, making it a \qbpn.

\subparagraph{Every pn has a polarized core.}

The encoding of \QBNs into polarized \good\ \qpns is sound and complete.
In fact, we can associate a \QBN\ to any \good \qpn (first reducing it to its normal form).

\begin{proposition}[Polarized core of a Normal Form]
Every \emph{normal} \pn $\somepn \der F$ consists of a polarized \pn $\pol{\somepn}$ on top of the formula tree of $F$.
\end{proposition}
This generalizes in the obvious way to $\somepn \der \Delta$.
For example, consider the proof-nets of \cref{fig:modularity_pn}: each has a formula tree drawn with \textcolor{red}{dashed} edges below its polarized core with \textcolor{black}{solid} edges.

Putting everything together, we can associate a \qqbn\ to (the normal form of) any \good \qbpn.
In the next section, we prove they have the \emph{same semantics} (\Cref{th:sem_coherence}).

\subsection{Compositionality and Modularity}
\label{sec:PNcompositionality}

We have now built all the ingredients to satisfy the desiderata outlined in the introduction.

\subparagraph{Denotational Semantics.}

We denote by $\Nm{\somepn}$ the set of all names appearing in $\somepn$.
\begin{definition}[Semantics of a \qpn]
\label{def:sem}
Let $(\somepn,\Phi)$ be a \qbpn.
Its semantics is:
\begin{equation*}
\den{\somepn} \defeq \sum_{\Nm{\somepn}\difset\Nm{\Delta}} \left(\bigodot_{\bbox\in \Boxes{\somepn}} \Phi(\bbox)\right)
\quad\text{where $\Delta$ is the conclusions of $\somepn$.}
\end{equation*}
\end{definition}

\begin{remark}\label{ex:var_ax}
If $\somepn$ is such that $\Boxes{\somepn}=\emptyset$, then $\den{\somepn}$ is the trivial factor $\ftone$.
\end{remark}

The invariance of the semantics via reduction is immediate, because both the boxes and the conclusions are invariants of the reduction.

\begin{proposition}[Invariance]
Let $\somepn$ be a \qpn.
\begin{enumerate}
\item
If $\somepn \Red \somepn'$ then $\somepn'$ is a \qpn with $\den{\somepn}=\den{\somepn'}$.
\item
Let $\somepn'$ be the normal form of $\somepn$, then $\den{\somepn} = \den{\somepn'}= \den {\pol{\somepn'}}$.
\end{enumerate}
\end{proposition}

The crucial question is \emph{compositionality}.
We prove that the semantics of a \qbpn $\somepn$ can be defined compositionally for any decomposition of $\somepn$ in sub-graphs.

\begin{theorem}[Compositionality]
\label{th:compositional}
Let $(\somepn,\Phi)$ be a \qpn, decomposed into two sub-graphs $\somepn_1$ and $\somepn_2$ with the obvious \qCPTs assignments given by restriction of $\Phi$.
Then:
\begin{equation*}
\den{\somepn} = \sum_{(\Nm{\Delta_1} \cup \Nm{\Delta_2})\difset\Nm{\Delta}} \big(\den{\somepn_1} \odot \den{\somepn_2}\big)
\end{equation*}
where $\Delta$ (\resp\ $\Delta_1$, $\Delta_2$) is the conclusions of $\somepn$ (\resp\ $\somepn_1$, $\somepn_2$).
\end{theorem}
The proof is similar to~\cite{BPN}, adapting to a graph setting the argument from~\cite{popl24}.

\subparagraph{Soundness and Completeness.}

In \Cref{sec:QBN_qpn}, we proved that the formalism of quantum proof-nets is sound and complete \wrt\ QBNs.
The semantics of a \qbpn\ and of the associated \qqbn\ are the same, essentially by definition.

\begin{theorem}[Completeness and Soundness with respect to QBNs]
\label{th:sem_coherence}
\hfill
\begin{itemize}
\item
To every \qqbn\ $\bn$ is associated a \good \qbpn $\somepn_\bn$ with $\sem{\bn} = \den{\somepn_{\bn}}$.
\item
To every \emph{\good} \qbpn\ $\somepn\der \Delta$ is associated a \qqbn\ $\bn_{\somepn}$ with $\den{\somepn} = \sum_{\Nm{\somepn} \difset \Nm{\Delta}}\sem {\bn_\somepn}$.
\end{itemize}
\end{theorem}

%\subparagraph{Observed and unobserved variables.}
%
%A fine point in the correspondence between \qqbn's and \pns deserves discussion.
%A \good\ \qbpn\ $\somepn\der\Delta$ encodes both a \QBN\ (the probabilistic model) and a \emph{query} to the model, expressed by the conclusions $\Delta$.
%As was formalized in \Cref{sec:QBN_qpn}, $\somepn\der\Delta$ corresponds to a \BN $\bn_\somepn$ over $\Nm{\somepn}$.
%While $\sem{\bn_\somepn}$ is a probability distribution over $\Nm{\somepn}$, $\den{\somepn}$ is the \emph{marginal distribution} of $\sem{\bn_\somepn}$ over $\Nm{\Delta} \subseteq \Nm{\somepn}$ (see \cref{th:sem_coherence}).
%A \good \pn $\somepn:\Delta$ defines a marginal distribution over $\Nm{\Delta}$.

\subparagraph{On Types and Modularity.}

The typing discipline enables modularity: two \qbpns of dual type $F$ and $F\orth$ are \emph{guaranteed} to compose well.
More generally, composition is defined by the standard $\cut$-rule of sequent calculus---we say that $\R_1,\R_2$ below have \definitive{compatibles types}:
\begin{equation*}
\infer[\cut]{\somepn\der \Gamma,\Delta}{\somepn_1 \der \Gamma, F & \somepn_2 \der \Delta, F\orth}
\end{equation*}
where $\somepn$ is obtained by connecting the conclusions (labelled by) $F$ and $F\orth$ with a $\cut$-node.

\begin{proposition}
\label{th:types}
If the \qbpns $\somepn_1$ and $\somepn_2$ have compatible types, their composition is a \qbpn.
Plugging together the DAGs respectively associated to $\somepn_1$ and $\somepn_2$ produces a DAG.
\end{proposition}

\cref{fig:modularity_pn} revisits the motivational example of \Cref{fig:modular} in our typed setting. \cref{fig:modularity_pn_bad} presents a non-example of typing for $\N_3$ in \Cref{fig:modular}: this graph is not a proof-net due to the highlighted cycle.
Types enable composition, by acting as an interface which encodes (via $\otimes$ and $\parr$) the DAG's structure.

\begin{figure}
\begin{minipage}[b]{0.60\linewidth}
\begin{adjustbox}{}
\begin{tikzpicture}[baseline=(base)]
\coordinate (base) at (0,-2.6);
\node at (1.5,0) {$\somepn_0$};
\begin{scope}[every node/.style={rectangle,draw=none}, every path/.style={draw=black}]
	\node[draw=black,minimum width=3em] (bB) at (0,0) {$B$};
	\node[draw=black,minimum width=3em] (bC) at (3,0) {$C$};

	\node (cB) at (1.5,-.75) {$\cut$};
	\path (0.3,-.25) |- node[named edge,pos=.5,above right]{$B^+$} (cB) -| node[named edge,pos=.5,above left]{$B^-$} (2.6,-.25);
	\path (-.2,-.25) -- node[named edge,pos=.75,left]{$A^-$} ++(0,-.5);
	\path (3.3,-.25) -- node[named edge,pos=.75,right]{$C^+$} ++(0,-.5);
\end{scope}
\begin{scope}[every node/.style={rectangle,draw=none}, every path/.style={draw=red,dashed}]
	\node (parr) at (1.5,-1.5) {\color{red}$\parr$};
	\path[out=-90,in=180] (-.2,-.8) edge (parr);
	\path[out=-90,in=0] (3.3,-.8) edge (parr);
	\path (parr) -- node[named edge,pos=.75]{\color{red}$A^+\lollipop C^+$} ++(0,-.5);
\end{scope}
\end{tikzpicture}
\qquad
\begin{tikzpicture}
\node at (6.5,-.5) {$\somepn_1$};
\begin{scope}[every node/.style={rectangle,draw=none}, every path/.style={draw=black}]
	\node[draw=black,minimum width=3em] (bE) at (5,-.5) {$E$};
	\node[draw=black,minimum width=3em] (bD) at (10,-.5) {$D$};
	\node[draw=black,minimum width=3em] (bA) at (8,-.5) {$A$};

	\node (cE) at (6.5,-1.5) {$\cut$};
	\node (pE) at (8.9,-1) {$\cc$};
	\path (5.3,-.75) |- node[named edge,pos=.5,above right]{$E^+$} (cE) -| node[named edge,pos=.2,above left]{$E^-$} (pE);
	\path[out=180,in=-90] (pE) edge (7.8,-.75);
	\path[out=0,in=-90] (pE) edge (9.6,-.75);
	\path (9.8,-.75) -- node[named edge,pos=1,left]{$C^-$} ++(0,-1.25);
	\path (10.3,-.75) -- node[named edge,pos=.9]{$D^+$} ++(0,-2.25);
	\path (8.3,-.75) -- node[named edge,pos=1]{$A^+$} ++(0,-1.25);
\end{scope}
\begin{scope}[every node/.style={rectangle,draw=none}, every path/.style={draw=red,dashed}]
	\node (parr) at (9,-2.5) {\color{red}$\tens$};
	\path[out=-90,in=180] (8.3,-2) edge (parr);
	\path[out=-90,in=0] (9.8,-2) edge (parr);
	\path (parr) -- node[named edge,pos=.75,left]{\color{red}$(A^+\lollipop C^+)\orth$} ++(0,-.5);
\end{scope}
\end{tikzpicture}
\end{adjustbox}
\caption{Modularity through typing.}
\label{fig:modularity_pn}
\end{minipage}
\hfill
\begin{minipage}[b]{0.38\linewidth}
\begin{adjustbox}{}
\begin{tikzpicture}
\node at (6.5,-6.5) {$\somepn_2$};
\begin{scope}[every node/.style={rectangle,draw=none}, every path/.style={draw=black}]
	\node[draw=black,minimum width=3em] (bE) at (5,-6.5) {$E$};
	\node[draw=black,minimum width=3em] (bD) at (8,-6.5) {$D$};
	\node[draw=black,minimum width=3em] (bA) at (10,-6.5) {$A$};

	\node (cE) at (6.5,-7.5) {$\cut$};
	\node (pE) at (8.9,-7) {$\cc$};
	\node (cD) at (9,-8) {$\cut$};
	\node (pD) at (10,-7.5) {$\cc$};
	\node (axD) at (11,-7) {$\ax$};
\end{scope}
\begin{scope}[every node/.style={rectangle,draw=none}, every path/.style={draw=red,dashed}]
	\node (parr) at (9,-8.5) {\color{red}$\tens$};
\end{scope}
\begin{scope}[every path/.style={thick,draw=yellow,-,line width=4pt}]
	\path[out=180,in=-90] (pD) edge (9.8,-6.75);
	\path (8.3,-6.75) |- (cD) -| (pD);
	\path (7.6,-6.75) -- ++(0,-1.25);
	\path (10.3,-6.75) -- ++(0,-1.25);
	\path[out=-90,in=180] (7.6,-8) edge (parr);
	\path[out=-90,in=0] (10.3,-8) edge (parr);
\end{scope}
\begin{scope}[every node/.style={rectangle,draw=none}, every path/.style={draw=black}]
	\path (5.3,-6.75) |- node[named edge,pos=.5,above right]{$E^+$} (cE) -| node[named edge,pos=.2,above left]{$E^-$} (pE);
	\path[out=180,in=-90] (pE) edge (7.8,-6.75);
	\path[out=0,in=-90] (pE) edge (9.6,-6.75);
	\path (8.3,-6.75) |- node[named edge,pos=.5,above right]{$D^+$} (cD) -| node[named edge,pos=.5,above left]{$D^-$} (pD);
	\path[out=180,in=-90] (pD) edge (9.8,-6.75);
	\path[out=0,in=-90] (pD) edge (10.5,-7);
	\path (axD) -| (10.5,-7);
	\path (axD) -| node[named edge,pos=.95]{$D^+$} (11.5,-9);
	\path (7.6,-6.75) -- node[named edge,pos=1,left]{$C^-$} ++(0,-1.25);
	\path (10.3,-6.75) -- node[named edge,pos=1,right]{$A^+$} ++(0,-1.25);
\end{scope}
\begin{scope}[every node/.style={rectangle,draw=none}, every path/.style={draw=red,dashed}]
	\path[out=-90,in=180] (7.6,-8) edge (parr);
	\path[out=-90,in=0] (10.3,-8) edge (parr);
	\path (parr) -- node[named edge,pos=.75,left]{\color{red}$(A^+\lollipop C^+)\orth$} ++(0,-.5);
\end{scope}
\end{tikzpicture}
\end{adjustbox}
\caption{Non-example of proof-net.}
\label{fig:modularity_pn_bad}
\end{minipage}
\end{figure}
